%% LyX 1.1 created this file.  For more info, see http://www.lyx.org/.
%% Do not edit unless you really know what you are doing.
\documentclass[a4paper,english]{article}
\usepackage[T1]{fontenc}
\usepackage[latin1]{inputenc}
\usepackage{babel}
\IfFileExists{url.sty}{\usepackage{url}}
                      {\newcommand{\url}{\texttt}}

\makeatletter

%%%%%%%%%%%%%%%%%%%%%%%%%%%%%% LyX specific LaTeX commands.
\providecommand{\LyX}{L\kern-.1667em\lower.25em\hbox{Y}\kern-.125emX\@}

%%%%%%%%%%%%%%%%%%%%%%%%%%%%%% User specified LaTeX commands.
\usepackage{neilbib}
\bibliographystyle{neilbib}
\title{Network Performance Analysis COM3270/6131}
\date{September 31st, 2001}


\author{Neil Lawrence\\{\tt neil@dcs.shef.ac.uk}}

\makeatother
\begin{document}

\title{Network Performance Analysis COM6131/COM3270\\
\url{http://www.dcs.shef.ac.uk/~neil/NPA/index.html}}


\date{October 1st, 2002}

\maketitle

\section*{\protect\centering Times}

\begin{description}
\item [Lectures]Tuesdays 2:10 pm - 16:00 pm (St Georges LT7) 
\item [Tutorials]Thursdays 12:10 pm - 13:00 pm (ARTS Tower LT3). First one
Week 3 - Tonatiuh Pena Centeno is your teaching assistant. 
\item [Evaluation]70\% Exam 30\% course work. 
\end{description}

\section*{\protect\centering Reading List}

The course is based primarily on \cite{Bertsekas:data92}. The relevant
parts of the book are as follows. 

\begin{description}
\item [Chapter~1]For background reading and review. 
\item [Chapter~3]Section 3.1 -- 3.5 For queueing models. 
\item [Chapter~4]In less detail: Section 4.1 -- 4.4, Section 4.5.2, 4.5.3. 
\end{description}
For additional background reading you may wish to consult some of
the books listed below, \cite{Tanenbaum:computer96} is the standard
text for Introductory courses, but has little coverage on queueing
theory. \cite{Chan:performance00} is an alternative text for details
on queues.

\bibliographystyle{neilbib}
\bibliography{lawrence_books}

\end{document}
